We will see later on that we will want to be able to identify when a certain token appeared in the tape.
Also, we would like to know when a certain position is reached or if a certain vector corresponds to a specific position.
This last task is not naturally done because of the extreme parallelism in the model.

For flagging a token we use the token embedding, and in the same way, for flagging a position we use the position encoding.
Either way, the strategy is exactly the same.

For each flag that we are looking to add, we will add one to the dimension embedding.
Notice that if we are working with a constant number of flags, the $\CoTtndn$ class does not change. {\color{orange} Not sure about this last sentence. Maybe I should expand the idea if I want to say it}

The idea for flagging is straight forward.
We will have an extra coordinate that will work as the flag.
This coordinate will have two possible values to be defined.
It will have the $\ON$ value $f_{\ON}$ if we are flagging its vector or the $\OFF$ value $f_{\OFF}$, otherwise.

For example, let's be $\TF$ a transformer and $\theta_{TE}$ its token embedding function.
We will define a transformer $\TF'$ with the same behavior, but it will have flagged every vector that came from the token $\sigma_f$.

For the sake of clarity, we will add the flag in the last coordinate of the embedding, but this can be done in any position.

Let's define each component of $\TF'$ based on the components of $\TF$.
First, let's see the token embedding function.
It will be the same except that in the new coordinate it will put a $V_{\ON}$ if the token that is consuming is the one to be flagged and a $0$ otherwise.

\begin{equation*}
    (\theta'_{TE}(\sigma))_j = \begin{cases}
        (\theta_{TE}(\sigma))_j  & \text{ if } j \neq d+1 \\
        f_{\ON}                  & \text{ if } j = d+1  \text{ and } \sigma = \sigma_f \\
        f_{\OFF}                 & \text{ if } j = d+1  \text{ and } \sigma \neq \sigma_f
    \end{cases}
\end{equation*}


For the position encoding is enough to always put a $0$ in the new coordinate.
Then is only necessary to extend the matrix of the other layers to the new dimension.
This can be done with a new row or column with $0$s.
The behavior of $\TF'$ will be the same as $\TF$.

This mechanism clearly works for adding more dimensions to the embedding.
We just add new flags with the same value for $f_{\ON}$ and $f_{\OFF}$.
